% =============================================================
%  Two-port networks — the general theory
%  Placed before the AC chapter: the line is one application of
%  this algebra, not its origin.  Section 4 (AC Power
%  Transmission) refers back here.
%  Figures are inline TikZ in the report colours.
% =============================================================
\section{Two-Port Networks}
\label{sec:two-port}

This appendix is background. Section~\ref{sec:ac-chapter} is self-contained
and nothing here is needed to follow it; what follows is for the reader who
wants to know where its equations come from and how far they reach.

The chain of matrices that Section~\ref{sec:ac-chapter} uses to describe a
transmission line is not a power-transmission construction. It is the
standard description of any linear network reached through two terminal
pairs, and it was in use for telephone repeaters and filters before it was
applied to power lines. Setting it out on its own terms shows the line to be
one instance of it rather than a special theory, and shows which of its
results are properties of a line and which are properties of any two-port at
all.

Two properties make the description worth the abstraction. The first is that
it hides everything inside the network: two ports and four complex numbers
describe the terminal behaviour completely, whatever the internal
arrangement. The second is that networks in cascade multiply, so a chain of
elements of any length collapses to one matrix of the same size.

\subsection{The port concept}
\label{sec:tp-ports}

A \emph{port} is a pair of terminals through which the current entering one
terminal equals the current leaving the other. That equality is the port
condition, and it is an assumption about how the network is connected, not a
property of the network itself: it fails if a third path --- a stray
capacitance to earth, a shared return, a ground loop --- carries current
between the two sides. A two-port is a network with exactly two such pairs,
Figure~\ref{fig:tp-generic}.

\begin{figure}[!ht]
  \centering
  \begin{tikzpicture}[
      x=0.62cm, y=0.62cm, line width=0.8pt,
      every node/.style={inner sep=1.5pt, font=\small},
      rail/.style={atFg, line width=0.8pt},
      nbox/.style={draw=atFg, fill=atBgSubtle, line width=0.8pt,
                   minimum width=2.108cm, minimum height=1.984cm},
      term/.style={circle, draw=atFg, fill=atBg, line width=0.8pt,
                   minimum size=4.2pt, inner sep=0pt},
      flow/.style={atFg, -{Triangle[length=2.4mm,width=2.0mm]},
                   line width=0.8pt},
    ]
    \node[nbox] at (5,1.6) {$N$};
    % rails
    \draw[rail] (0,3.2) -- (3.3,3.2);   \draw[rail] (6.7,3.2) -- (10,3.2);
    \draw[rail] (0,0)   -- (3.3,0);     \draw[rail] (6.7,0)   -- (10,0);
    % terminals
    \node[term] at (0,3.2) {};  \node[term] at (0,0) {};
    \node[term] at (10,3.2) {}; \node[term] at (10,0) {};
    % port currents: both drawn into the upper terminal
    \draw[flow] (-1.8,3.2) -- (-0.9,3.2);
    \node[left] at (-2.0,3.2) {$\underline{I}_1$};
    \draw[flow] (11.8,3.2) -- (10.9,3.2);
    \node[right] at (12.0,3.2) {$\underline{I}_2$};
    % port voltages
    \draw[flow] (0,2.8) -- (0,0.4);
    \node[right=2pt] at (0,1.6) {$\underline{U}_1$};
    \draw[flow] (10,2.8) -- (10,0.4);
    \node[left=2pt] at (10,1.6) {$\underline{U}_2$};
    % port brackets
    \node[font=\scriptsize, atFgMuted] at (0.1,-1.1) {port 1};
    \node[font=\scriptsize, atFgMuted] at (9.9,-1.1) {port 2};
  \end{tikzpicture}
  \caption{A two-port and its terminal quantities, drawn in the convention
           in which both port currents enter the network}
  \label{fig:tp-generic}
\end{figure}

Four quantities appear at the terminals: $\underline{U}_1$,
$\underline{I}_1$, $\underline{U}_2$ and $\underline{I}_2$. The network
imposes two linear relations between them, so any two may be chosen as
independent and the other two follow. Which two are chosen is the only
difference between the parameter sets below.

\subsection{Parameter sets}
\label{sec:tp-parameters}

Table~\ref{tab:tp-sets} lists the six choices. The impedance, admittance and
the two hybrid sets are drawn in the convention of
Figure~\ref{fig:tp-generic}, with both currents entering the network. The
chain sets reverse the reference direction at port 2, so that
$\underline{I}_2$ leaves the network towards whatever is connected next ---
without that reversal, cascading would alternate in sign at every junction.
Comparing results from two sources that differ only in this convention is
the most common error in using them.

\begin{table}[!ht]
  \centering
  \begin{tabularx}{\linewidth}{@{}l l X@{}}
    \toprule
    \textbf{Set} & \multicolumn{1}{c}{\textbf{Relation}} & \textbf{Typical use} \\
    \midrule
    Impedance $\underline{Z}$
      & $\underline{U} = \underline{Z} \cdot \underline{I}$
      & Networks joined in series; open-circuit measurements \\
    Admittance $\underline{Y}$
      & $\underline{I} = \underline{Y} \cdot \underline{U}$
      & Networks joined in parallel; nodal analysis \\
    Hybrid $\underline{h}$
      & $\left(\underline{U}_1, \underline{I}_2\right)$ from
        $\left(\underline{I}_1, \underline{U}_2\right)$
      & Amplifying devices, where the input is a current and the output a
        voltage \\
    Inverse hybrid $\underline{g}$
      & $\left(\underline{I}_1, \underline{U}_2\right)$ from
        $\left(\underline{U}_1, \underline{I}_2\right)$
      & The dual of the above \\
    Chain $\underline{T}$ (ABCD)
      & port 1 from port 2
      & Elements in cascade; transmission paths \\
    Inverse chain
      & port 2 from port 1
      & The same, traversed in the other direction \\
    \bottomrule
  \end{tabularx}
  \caption{The six parameter sets of a two-port, distinguished only by which
           two terminal quantities are taken as independent~[19], [20]}
  \label{tab:tp-sets}
\end{table}

The sets are equivalent wherever they all exist, and each is undefined for
some networks: an ideal transformer has no impedance matrix, an ideal
series element no admittance matrix. The chain matrix exists unless the
network cannot transfer anything between the ports at all.

\subsection{Chain parameters and cascading}
\label{sec:tp-chain}

The chain form writes port 1 in terms of port 2,
\begin{equation}
  \begin{bmatrix} \underline{U}_1 \\ \underline{I}_1 \end{bmatrix}
  =
  \underline{T} \cdot
  \begin{bmatrix} \underline{U}_2 \\ \underline{I}_2 \end{bmatrix},
  \qquad
  \underline{T} =
  \begin{bmatrix} \underline{A} & \underline{B} \\
                  \underline{C} & \underline{D} \end{bmatrix}
  \label{eq:tp-chain}
\end{equation}
with $\underline{I}_2$ leaving the network. The four entries have a meaning
one at a time. With port 2 open, $\underline{I}_2 = 0$ and $\underline{A} =
\underline{U}_1/\underline{U}_2$ is the reciprocal voltage ratio, while
$\underline{C} = \underline{I}_1/\underline{U}_2$ is a transfer admittance.
With port 2 shorted, $\underline{U}_2 = 0$ and $\underline{B} =
\underline{U}_1/\underline{I}_2$ is a transfer impedance, $\underline{D} =
\underline{I}_1/\underline{I}_2$ the reciprocal current ratio. So
$\underline{A}$ and $\underline{D}$ are dimensionless, $\underline{B}$ is an
impedance and $\underline{C}$ an admittance --- which is why the four cannot
be compared with each other as magnitudes.

The reason this set is the one used for transmission is cascading. If two
two-ports are connected so that port 2 of the first is port 1 of the second,
Figure~\ref{fig:tp-cascade}, the intermediate quantities are shared and
eliminate, leaving
\begin{equation}
  \underline{T} = \underline{T}_a \cdot \underline{T}_b,
  \qquad\text{and in general}\qquad
  \underline{T} = \prod_{i=1}^{n} \underline{T}_i
  \label{eq:tp-cascade}
\end{equation}
in the order the elements are traversed. The product does not commute:
reversing two elements of a chain generally changes the result, which is the
algebraic statement that it matters where in a path an element sits. A
compensating element at the middle of a line and the same element at its end
are different networks, and~\eqref{eq:tp-cascade} is what distinguishes
them.

\begin{figure}[!ht]
  \centering
  \begin{tikzpicture}[
      x=0.62cm, y=0.62cm, line width=0.8pt,
      every node/.style={inner sep=1.5pt, font=\small},
      rail/.style={atFg, line width=0.8pt},
      nbox/.style={draw=atFg, fill=atBgSubtle, line width=0.8pt,
                   minimum width=1.736cm, minimum height=1.86cm},
      term/.style={circle, draw=atFg, fill=atBg, line width=0.8pt,
                   minimum size=4.2pt, inner sep=0pt},
      flow/.style={atFg, -{Triangle[length=2.4mm,width=2.0mm]},
                   line width=0.8pt},
    ]
    \node[nbox] at (3.4,1.5) {$\underline{T}_a$};
    \node[nbox] at (9.6,1.5) {$\underline{T}_b$};
    % left block
    \draw[rail] (0,3.0) -- (2.0,3.0);  \draw[rail] (4.8,3.0) -- (8.2,3.0);
    \draw[rail] (0,0)   -- (2.0,0);    \draw[rail] (4.8,0)   -- (8.2,0);
    % right block
    \draw[rail] (11.0,3.0) -- (13.0,3.0);
    \draw[rail] (11.0,0)   -- (13.0,0);
    % intermediate port terminals
    \node[term] at (6.5,3.0) {};  \node[term] at (6.5,0) {};
    % terminals
    \node[term] at (0,3.0) {};   \node[term] at (0,0) {};
    \node[term] at (13,3.0) {};  \node[term] at (13,0) {};
    % currents
    \draw[flow] (-1.7,3.0) -- (-0.9,3.0);
    \node[left] at (-1.9,3.0) {$\underline{I}_1$};
    \draw[flow] (13.9,3.0) -- (14.7,3.0);
    \node[right] at (14.9,3.0) {$\underline{I}_2$};
    \draw[flow] (5.1,3.0) -- (5.9,3.0);
    \node[above=1pt] at (5.5,3.1) {$\underline{I}_m$};
    % voltages
    \draw[flow] (0,2.6) -- (0,0.4);
    \node[right=2pt] at (0,1.5) {$\underline{U}_1$};
    \draw[flow] (6.5,2.6) -- (6.5,0.4);
    \node[right=2pt] at (6.5,1.5) {$\underline{U}_m$};
    \draw[flow] (13,2.6) -- (13,0.4);
    \node[left=2pt] at (13,1.5) {$\underline{U}_2$};
  \end{tikzpicture}
  \caption{Two two-ports in cascade: the intermediate quantities are common
           to both and eliminate, so the chain matrices multiply}
  \label{fig:tp-cascade}
\end{figure}

\subsection{Other interconnections}
\label{sec:tp-interconnection}

Cascade is one of five ways two two-ports can be joined, and each has its
own parameter set in which the combination is a simple sum,
Table~\ref{tab:tp-interconnection}. The pattern is the same throughout: pick
the set whose independent variables are the ones the connection forces to be
shared, and the matrices add.

\begin{table}[!ht]
  \centering
  \begin{tabularx}{\linewidth}{@{}l l l X@{}}
    \toprule
    \textbf{Connection} & \multicolumn{1}{c}{\textbf{Shared}}
      & \multicolumn{1}{c}{\textbf{Rule}} & \textbf{Where it is met} \\
    \midrule
    Cascade & port 2 of one is port 1 of the next
      & $\underline{T} = \underline{T}_a \cdot \underline{T}_b$
      & Line sections, a series bank between two half-lines, a transformer
        ahead of a line \\
    Series & both currents
      & $\underline{Z} = \underline{Z}_a + \underline{Z}_b$
      & Windings in series, impedances in a common current path \\
    Parallel & both voltages
      & $\underline{Y} = \underline{Y}_a + \underline{Y}_b$
      & Parallel circuits of a double-circuit line, transformers in
        parallel \\
    Series-parallel & $\underline{I}_1$ and $\underline{U}_2$
      & $\underline{h} = \underline{h}_a + \underline{h}_b$
      & Feedback arrangements in amplifiers \\
    Parallel-series & $\underline{U}_1$ and $\underline{I}_2$
      & $\underline{g} = \underline{g}_a + \underline{g}_b$
      & The dual of the above \\
    \bottomrule
  \end{tabularx}
  \caption{The five interconnections of two two-ports and the parameter set
           in which each is additive~[19], [20]}
  \label{tab:tp-interconnection}
\end{table}

Cascade is the odd one out, and the difference is worth stating plainly
because the word ``series'' is used loosely for both. In a cascade the
networks are passed through one after another and the matrices
\emph{multiply}; in a series connection they sit side by side with their
ports stacked, both carry the same currents, and the matrices \emph{add},
Figure~\ref{fig:tp-connections}.
A capacitor bank inserted in a line is a cascade, not a series connection of
two-ports, even though the capacitor is in series with the conductor: what
is in series there is a circuit element within one port, not two ports with
each other.

\begin{figure}[!ht]
  \centering
  \begin{tikzpicture}[
      x=0.52cm, y=0.52cm, line width=0.8pt,
      every node/.style={inner sep=1.5pt, font=\small},
      rail/.style={atFg, line width=0.8pt},
      wbox/.style={draw=atFg, fill=atBgSubtle, line width=0.8pt,
                   minimum width=1.25cm, minimum height=2.184cm},
      sbox/.style={draw=atFg, fill=atBgSubtle, line width=0.8pt,
                   minimum width=1.25cm, minimum height=1.664cm},
      term/.style={circle, draw=atFg, fill=atBg, line width=0.8pt,
                   minimum size=3.6pt, inner sep=0pt},
      plab/.style={font=\scriptsize, atFgMuted},
    ]
    % ---------------- (a) cascade -------------------------------
    \begin{scope}[shift={(0,0)}]
      \draw[rail] (0,2.1) -- (1.9,2.1);   \draw[rail] (4.3,2.1) -- (7.7,2.1);
      \draw[rail] (0,-2.1) -- (1.9,-2.1); \draw[rail] (4.3,-2.1) -- (7.7,-2.1);
      \draw[rail] (10.1,2.1) -- (12,2.1); \draw[rail] (10.1,-2.1) -- (12,-2.1);
      \node[wbox] at (3.1,0) {$a$};
      \node[wbox] at (8.9,0) {$b$};
      \node[term] at (0,2.1) {};  \node[term] at (0,-2.1) {};
      \node[term] at (12,2.1) {}; \node[term] at (12,-2.1) {};
      \node[plab] at (6,-4.4)
        {(a) cascade: $\underline{T} = \underline{T}_a \cdot \underline{T}_b$};
    \end{scope}
    % ---------------- (b) series --------------------------------
    \begin{scope}[shift={(15.5,0)}]
      \draw[rail] (0,4.0) -- (3.8,4.0);   \draw[rail] (6.2,4.0) -- (10,4.0);
      \draw[rail] (3.8,0.8) -- (0,0.8) -- (0,-0.8) -- (3.8,-0.8);
      \draw[rail] (6.2,0.8) -- (10,0.8) -- (10,-0.8) -- (6.2,-0.8);
      \draw[rail] (0,-4.0) -- (3.8,-4.0); \draw[rail] (6.2,-4.0) -- (10,-4.0);
      \node[sbox] at (5,2.4) {$a$};
      \node[sbox] at (5,-2.4) {$b$};
      \node[term] at (0,4.0) {};  \node[term] at (0,-4.0) {};
      \node[term] at (10,4.0) {}; \node[term] at (10,-4.0) {};
      \node[plab] at (5,-6.0)
        {(b) series: $\underline{Z} = \underline{Z}_a + \underline{Z}_b$};
    \end{scope}
  \end{tikzpicture}
  \caption{Cascade against series connection. In the cascade the networks
           are passed through one after the other and the chain matrices
           multiply; in the series connection the ports are stacked, both
           networks carry the same currents, and the impedance matrices add}
  \label{fig:tp-connections}
\end{figure}

The addition rules carry a condition that the cascade does not. Stacking two
networks can destroy the port condition inside them: current that entered
one terminal of a sub-network may return through the other sub-network
rather than through its partner terminal, and then the sub-network no longer
has the matrix it was given alone, so the sum is wrong. Whether a particular
pairing is safe is settled by a test due to Brune, applied to the two
networks as connected~[20]. Two cases are safe without testing: either
sub-network having an isolating transformer at one port, and --- in power
work the common case --- two networks that share a common return and are
each symmetric about it. The cascade needs no such test, because it adds no
new connection between terminals that were not already a port.

The parallel rule is the one met most often in transmission. Two identical
circuits on the same route, each with parameters $\underline{A} \ldots
\underline{D}$, give
\begin{equation}
  \underline{A}_{\parallel} = \underline{A},
  \qquad
  \underline{B}_{\parallel} = \frac{\underline{B}}{2},
  \qquad
  \underline{C}_{\parallel} = 2 \cdot \underline{C},
  \qquad
  \underline{D}_{\parallel} = \underline{D}
  \label{eq:tp-parallel}
\end{equation}
which is the statement that two circuits halve the transfer impedance and so
double the maximum transfer, while the voltage ratio is unchanged. Losing
one of them restores $\underline{B}$, which is the contingency case of
Section~\ref{sec:as-contingency}. Equation~\eqref{eq:tp-parallel} follows
from the admittance rule, and is quicker to use than converting twice.

\subsection{Elementary two-ports}
\label{sec:tp-elements}

Every ladder network is built from two elements: an impedance in series with
one rail and an admittance across the rails, Figure~\ref{fig:tp-elements}(a)
and~(b). Their chain matrices are
\begin{equation}
  \underline{T}_{series} =
  \begin{bmatrix} 1 & \underline{Z} \\ 0 & 1 \end{bmatrix},
  \qquad
  \underline{T}_{shunt} =
  \begin{bmatrix} 1 & 0 \\ \underline{Y} & 1 \end{bmatrix}
  \label{eq:tp-elements}
\end{equation}
and everything else in Table~\ref{tab:tp-matrices} is a product of these
two, taken in the order the elements are met.

\begin{figure}[!ht]
  \centering
  \begin{tikzpicture}[
      x=0.62cm, y=0.62cm, line width=0.8pt,
      every node/.style={inner sep=1.5pt, font=\small},
      rail/.style={atFg, line width=0.8pt},
      box/.style={draw=atFg, fill=atBgSubtle, line width=0.8pt,
                  minimum width=1.0cm, minimum height=0.42cm},
      vbox/.style={draw=atFg, fill=atBgSubtle, line width=0.8pt,
                   minimum width=0.5cm, minimum height=1.0cm},
      term/.style={circle, draw=atFg, fill=atBg, line width=0.8pt,
                   minimum size=3.6pt, inner sep=0pt},
      plab/.style={font=\scriptsize, atFgMuted},
    ]
    % ---- (a) series impedance ----------------------------------
    \begin{scope}[shift={(0,0)}]
      \draw[rail] (0,2.6) -- (1.79,2.6);
      \node[box] at (2.6,2.6) {$\underline{Z}$};
      \draw[rail] (3.41,2.6) -- (5.2,2.6);
      \draw[rail] (0,0) -- (5.2,0);
      \node[term] at (0,2.6) {}; \node[term] at (0,0) {};
      \node[term] at (5.2,2.6) {}; \node[term] at (5.2,0) {};
      \node[plab] at (2.6,-1.1) {(a) series impedance};
    \end{scope}
    % ---- (b) shunt admittance ----------------------------------
    \begin{scope}[shift={(7.8,0)}]
      \draw[rail] (0,2.6) -- (5.2,2.6);
      \draw[rail] (0,0) -- (5.2,0);
      \draw[rail] (2.6,2.6) -- (2.6,2.10);
      \node[vbox] at (2.6,1.3) {$\underline{Y}$};
      \draw[rail] (2.6,0.50) -- (2.6,0);
      \node[term] at (0,2.6) {}; \node[term] at (0,0) {};
      \node[term] at (5.2,2.6) {}; \node[term] at (5.2,0) {};
      \node[plab] at (2.6,-1.1) {(b) shunt admittance};
    \end{scope}
    % ---- (c) T-section -----------------------------------------
    \begin{scope}[shift={(0,-5.6)}]
      \draw[rail] (0,2.6) -- (0.59,2.6);
      \node[box] at (1.4,2.6) {$\underline{Z}_1$};
      \draw[rail] (2.21,2.6) -- (2.99,2.6);
      \node[box] at (3.8,2.6) {$\underline{Z}_2$};
      \draw[rail] (4.61,2.6) -- (5.2,2.6);
      \draw[rail] (0,0) -- (5.2,0);
      \draw[rail] (2.6,2.6) -- (2.6,2.10);
      \node[vbox] at (2.6,1.3) {$\underline{Y}$};
      \draw[rail] (2.6,0.50) -- (2.6,0);
      \node[term] at (0,2.6) {}; \node[term] at (0,0) {};
      \node[term] at (5.2,2.6) {}; \node[term] at (5.2,0) {};
      \node[plab] at (2.6,-1.1) {(c) T-section};
    \end{scope}
    % ---- (d) pi-section ----------------------------------------
    \begin{scope}[shift={(7.8,-5.6)}]
      \draw[rail] (0,2.6) -- (1.79,2.6);
      \node[box] at (2.6,2.6) {$\underline{Z}$};
      \draw[rail] (3.41,2.6) -- (5.2,2.6);
      \draw[rail] (0,0) -- (5.2,0);
      \draw[rail] (1.0,2.6) -- (1.0,2.10);
      \node[vbox] at (1.0,1.3) {$\underline{Y}_1$};
      \draw[rail] (1.0,0.50) -- (1.0,0);
      \draw[rail] (4.2,2.6) -- (4.2,2.10);
      \node[vbox] at (4.2,1.3) {$\underline{Y}_2$};
      \draw[rail] (4.2,0.50) -- (4.2,0);
      \node[term] at (0,2.6) {}; \node[term] at (0,0) {};
      \node[term] at (5.2,2.6) {}; \node[term] at (5.2,0) {};
      \node[plab] at (2.6,-1.1) {(d) $\pi$-section};
    \end{scope}
  \end{tikzpicture}
  \caption{The two elementary two-ports and the two ladder sections built
           from them}
  \label{fig:tp-elements}
\end{figure}

\begin{table}[!ht]
  \centering
  \begin{tabularx}{\linewidth}{@{}l l X@{}}
    \toprule
    \textbf{Network} & \multicolumn{1}{c}{$\underline{A}$ \ldots $\underline{D}$}
      & \textbf{Remark} \\
    \midrule
    Series $\underline{Z}$
      & $1,\; \underline{Z},\; 0,\; 1$
      & $\underline{A} = 1$: no shunt path, so the open-circuit voltage
        ratio is unity \\
    Shunt $\underline{Y}$
      & $1,\; 0,\; \underline{Y},\; 1$
      & $\underline{B} = 0$: the ports are directly connected \\
    Symmetric $\pi$
      & $1 + \frac{\underline{Z} \cdot \underline{Y}}{2},\;
         \underline{Z},\;
         \underline{Y} \cdot \left(1 + \frac{\underline{Z} \cdot
         \underline{Y}}{4}\right),\;
         1 + \frac{\underline{Z} \cdot \underline{Y}}{2}$
      & $\underline{Y}/2$ at each end; the nominal line section
        of~\eqref{eq:ac-abcd-pi} \\
    Ideal transformer
      & $n,\; 0,\; 0,\; 1/n$
      & $\underline{B} = \underline{C} = 0$: a pure scaling of both
        quantities \\
    Uniform line
      & $\cosh \underline{\theta},\;
         \underline{Z}_w \cdot \sinh \underline{\theta},\;
         \sinh \underline{\theta} / \underline{Z}_w,\;
         \cosh \underline{\theta}$
      & $\underline{\theta} = \underline{\gamma} \cdot l$; the limit of
        infinitely many $\pi$-sections \\
    \bottomrule
  \end{tabularx}
  \caption{Chain matrices of the elementary two-ports and of the uniform
           line~[19], [21]}
  \label{tab:tp-matrices}
\end{table}

The last row is the one that matters for a transmission line and is derived
in Section~\ref{sec:ac-wave}; it is included here to show that it is the
same object as the others, not a separate theory.

\subsection{Reciprocity, symmetry and losslessness}
\label{sec:tp-properties}

Three structural properties of a network show themselves directly in the
four entries.

\paragraph{Reciprocity.} A network built only from impedances, whatever
their arrangement, satisfies
\begin{equation}
  \underline{A} \cdot \underline{D} - \underline{B} \cdot \underline{C} = 1
  \label{eq:tp-reciprocity}
\end{equation}
Physically: a source at port 1 and an ammeter at port 2 read the same as the
source and the ammeter exchanged. It fails for networks containing an
independent source or a non-reciprocal element such as a gyrator or a
magnetised ferrite. Because the determinant is unity, the product
in~\eqref{eq:tp-cascade} preserves it, so any cascade of reciprocal
networks is reciprocal.

\paragraph{Symmetry.} If the network looks the same from both ports then
\begin{equation}
  \underline{A} = \underline{D}
  \label{eq:tp-symmetry}
\end{equation}
Every uniform line is symmetric; a $\pi$-section with unequal shunt branches
is not, and neither is a bank placed anywhere other than at the midpoint.
Symmetry is the condition under which the image impedance of
Section~\ref{sec:tp-terminated} is the same at both ends.

\paragraph{Losslessness.} A network of pure reactances has $\underline{A}$
and $\underline{D}$ real and $\underline{B}$ and $\underline{C}$ purely
imaginary. The departure of $\underline{A}$ from a real number is therefore
a direct measure of the loss in a line, and is why the lossless
approximations of Section~\ref{sec:ac-chapter} can be stated by dropping
real parts rather than by rebuilding the model.

\subsection{Terminating a two-port}
\label{sec:tp-terminated}

A two-port on its own is not a circuit. Closing port 2 with an impedance
$\underline{Z}_L$ makes $\underline{U}_2 = \underline{Z}_L \cdot
\underline{I}_2$, and~\eqref{eq:tp-chain} gives the impedance the source at
port 1 sees:
\begin{equation}
  \underline{Z}_{in}
    = \frac{\underline{U}_1}{\underline{I}_1}
    = \frac{\underline{A} \cdot \underline{Z}_L + \underline{B}}
           {\underline{C} \cdot \underline{Z}_L + \underline{D}}
  \label{eq:tp-zin}
\end{equation}
This is the single most used result of the whole description: every
statement about what a network does to the thing behind it is a special case
of it. With port 2 open it gives $\underline{A}/\underline{C}$, with port 2
shorted $\underline{B}/\underline{D}$.

One termination is distinguished. On a symmetric network there is a value of
$\underline{Z}_L$ for which $\underline{Z}_{in} = \underline{Z}_L$, so that
the network is invisible to its surroundings and may be repeated any number
of times without changing anything. Setting $\underline{Z}_{in} =
\underline{Z}_L$ in~\eqref{eq:tp-zin} with $\underline{A} = \underline{D}$
gives the \emph{image} or \emph{characteristic} impedance
\begin{equation}
  \underline{Z}_c = \sqrt{\frac{\underline{B}}{\underline{C}}},
  \qquad
  \underline{\theta} = \operatorname{arcosh} \underline{A}
  \label{eq:tp-image}
\end{equation}
with $\underline{\theta}$ the image transfer constant, the complex number by
which a quantity is scaled and shifted in passing through the network.
Substituting the uniform-line row of Table~\ref{tab:tp-matrices} returns
$\underline{Z}_c = \underline{Z}_w$ and $\underline{\theta} =
\underline{\gamma} \cdot l$: the wave impedance of a line and its electrical
length are not line-specific quantities but this general pair, evaluated for
that network. The natural power of Section~\ref{sec:natural-power} is the
same statement in power terms --- a line closed in its own image impedance.

\subsection{The same algebra in other domains}
\label{sec:tp-domains}

Nothing above used the fact that the quantities were electrical. The
derivation needs only a pair of variables at each port whose product is
power: one that is common to the elements in parallel and one that is common
to the elements in series. Table~\ref{tab:tp-domains} lists the pairs for
which the identical algebra is standard practice.

\begin{table}[!ht]
  \centering
  \begin{tabularx}{\linewidth}{@{}l l l X@{}}
    \toprule
    \textbf{Domain} & \multicolumn{1}{c}{\textbf{Across}}
      & \multicolumn{1}{c}{\textbf{Through}} & \textbf{Where the chain form is used} \\
    \midrule
    Electrical & Voltage & Current
      & Lines, filters, matching networks, machine equivalents \\
    Translational & Velocity & Force
      & Vibration isolators, drive trains, transducer chains \\
    Rotational & Angular velocity & Torque
      & Shaft lines, gearboxes, torsional models of a turbine set \\
    Acoustic & Sound pressure & Volume velocity
      & Ducts, horns, mufflers, loudspeaker enclosures~[22] \\
    Thermal & Temperature difference & Heat flow
      & Layered walls, heat sinks, cable thermal ratings \\
    \bottomrule
  \end{tabularx}
  \caption{Domains in which the two-port algebra is used unchanged, with the
           pair of port variables whose product is power~[20], [22]}
  \label{tab:tp-domains}
\end{table}

A transducer is simply a two-port whose ports lie in two different rows of
that table: a loudspeaker is an electrical port and an acoustic one, a motor
an electrical port and a rotational one, a piezoelectric stack an electrical
port and a translational one. Cascading them is the same matrix product, and
it is how a complete electro-mechanical chain is reduced to one matrix.

Two further uses are worth naming because they are reached from a different
starting point and arrive at the same matrix.

\paragraph{Ray and beam optics.} A sequence of lenses, mirrors and free
spaces is handled by a $2 \times 2$ \emph{ray transfer matrix} acting on the
ray height and its slope, multiplied in the order the light meets the
elements~[23]. A thin lens is $\left[1, 0; -1/f, 1\right]$ --- the shunt
element of~\eqref{eq:tp-elements} with $\underline{Y} = -1/f$ --- and a free
space of length $d$ is $\left[1, d; 0, 1\right]$, the series element. The
resonator condition that the beam reproduce itself after one round trip is
the image-impedance condition~\eqref{eq:tp-image} in another vocabulary.

\paragraph{High-frequency networks.} Above a few hundred megahertz a
terminal voltage is not measurable and the scattering parameters are used in
place of the chain matrix. They are a change of variables, not a different
theory: the scattering matrix is the same network expressed in incident and
reflected waves against a reference impedance, and converts back to
$\underline{A} \ldots \underline{D}$ in closed form. Cascading in that
vocabulary uses the transfer-scattering matrix, which multiplies exactly
as~\eqref{eq:tp-cascade} does.

\subsection{What the description assumes}
\label{sec:tp-limits}

The generality has a boundary, and the assumptions are easy to leave
unstated because they hold without comment in most of this report.

\begin{itemize}
  \setlength{\itemsep}{0.35em}
  \item \textbf{Linear and time-invariant.} The relations are linear with
        constant coefficients. A saturating transformer core, a switching
        converter or a thyristor-controlled bank is not described by one
        matrix; it is described by one matrix per operating point, or by a
        model outside this framework.
  \item \textbf{One frequency.} The entries are complex numbers at a stated
        frequency. A chain evaluated at 60\,Hz says nothing about its
        behaviour at 20\,Hz, which is why the sub-synchronous screening of
        Section~\ref{sec:sc-ssr} re-evaluates the same network over a range
        of frequencies rather than scaling one result.
  \item \textbf{The port condition holds.} Current entering a terminal
        leaves its partner. Earth-return paths and inter-circuit coupling
        break it, and a line with a significant zero-sequence path is a
        multi-port, not a two-port.
  \item \textbf{No internal sources.} Independent sources inside the network
        add a constant term that no $2 \times 2$ matrix carries. They are
        handled by superposition, or by moving them outside the two-port.
  \item \textbf{Balanced three-phase, for the power case.} The single-line
        two-port of Section~\ref{sec:ac-chapter} is the positive-sequence
        network. Unbalance requires the sequence networks separately.
\end{itemize}

Within those bounds the description is exact, and the exactness is the
point: the equations of Section~\ref{sec:ac-chapter} are not an engineering
approximation to a line but the general two-port result with the line's
own $\underline{A} \ldots \underline{D}$ substituted.

\FloatBarrier
